% Ad Atticum 12.12 --- headnote
% ad-atticum-12-12, c. 16 March 45 BC, Astura

\textit{Cicero to Atticus, written from Astura on the seventeenth
day before the Kalends of April --- Perseus: \textit{Asturae
xvii K.~Apr.~a.~709 (45)}. Tullia, Cicero's only daughter, had died
in mid-February of that year, perhaps a month before this letter.
Cicero, refusing to return to Rome, had withdrawn to Astura on the
Latian coast; from there he writes to Atticus almost every day,
not because he has anything to report but because the daily letter
is what is now holding him to ordinary life. The voice of this
letter is the voice of those months: spare, withdrawn, businesslike
on the surface, with the grief implied rather than spoken.}

\textit{Two matters in front of him. The first is the dowry: with
Tullia's marriage to Dolabella formally ended and Tullia herself
now dead, the money owed back to her estate has to be settled, and
Balbus's proposal to transfer the debt by assignment is on the
table. Cicero presses Atticus to close it. The second, and the
quiet centre of the letter, is the \textit{fanum} --- the shrine
he intends to build to consecrate Tullia. The island at Arpinum
could afford a true \textit{apoth\=eosis}, he says, but its very
remoteness might make the honour look smaller; his mind is now
turning to a site in the suburban \textit{horti} of Rome. The
second section, ostensibly about which philosophers should figure
in his dialogues, ends in the line that gives the letter its
weight: I have nothing to write to you about, but I write all the
same, to draw your letters out of me --- not that I expect anything
from them, but somehow I still expect.}
